\chapter{延伸阅读}
\markboth{延伸阅读}{延伸阅读}

\begin{ledetext}
这本书只画了地图，画不出地形。下面五本，每一本都比这本书更深，也更难。
不必按顺序读，挑一本你此刻最想弄明白的那件事。
\end{ledetext}

\begin{flowlist}
\flowitem{01}{《现代经济中的货币创造》，英格兰银行季刊，2014}{
本书第 7 章的直接依据。官方出版物，免费，二十页，把「贷款创造存款」这件事讲得比任何教科书都清楚。
如果你只读一份延伸材料，读这个。它会让你确信第 7 章不是奇谈怪论，而是央行自己写的常识。}
\flowitem{02}{《疯狂、恐慌与崩溃》，查尔斯·金德尔伯格}{
本书第 17 章那个五幕剧本的来源。它把四百年的泡沫史压缩成一个可复用的结构，
读完你会发现每一次新的狂热都在按同一个剧本演，只是换了道具。}
\flowitem{03}{《债务危机》，瑞·达利欧}{
第 18 章债务长周期的框架出处。这本书的价值在框架不在结论：
去杠杆的四种手段如何配比，决定了一个国家是走出 1930 年代还是走进日本的三十年。}
\flowitem{04}{《大衰退》，辜朝明}{
理解「资产负债表衰退」最好的入口，也是第 9 章利率传导为什么会失灵的完整版。
读完你会明白，为什么有时候把利率降到零，钱也不动。}
\flowitem{05}{《非对称风险》，纳西姆·塔勒布}{
第 22 章遍历性那一节的思想来源。它讲的其实只有一件事：
集合平均不等于时间平均，而你只有一条时间线。这句话比任何收益率都重要。}
\end{flowlist}

\bookbreak

\section{还想再往下走}

如果读完上面五本还想继续，往三个方向走。

\textbf{往制度走。}去读会计准则、破产法和清算制度。它们无聊到极点，
但第 25 章说过，信任的物理基础全在那里。理解了它们，你会对「一个国家的利率下限由它的法院决定」这句话有实感。

\textbf{往数据走。}把中国人民银行、美联储、欧洲央行和 IIF 的定期报告加进你的阅读列表。
它们是免费的一手材料，而绝大多数财经新闻只是对它们的二手转述，还经常转错。
自己看一遍原文，你会发现新闻标题和报告结论之间的距离有多大。

\textbf{往怀疑走。}找那些和本书结论相反的严肃论述来读。
比如关于货币数量论的辩护、关于被动投资并未损害价格发现的实证、关于私募信贷风险被高估的分析。
一个只读支持自己观点材料的人，读得越多越危险。

\cleardoublepage
